\chapter[Fundamentos teóricos]{Fundamentos teóricos e reprodução do TG}
\label{chap:fundamentos}

Este capítulo reconstrói a cinemática linear que sustenta a pesquisa.
O objetivo é obter a curvatura sem pressupor propagação nula, resolver
os vínculos dos modelos históricos e de Fierz--Pauli e estabelecer o
domínio das comparações com o TG \autocite{tg2003}. A contagem de
amplitudes será demonstrada por parametrizações e menores não nulos,
em vez de inferida apenas de exemplos numéricos. A reprodução das
expressões publicadas será distinguida de sua interpretação como
deformação observável e de sua utilização futura em uma PTA.

\section{Convenções e domínio da derivação}

Adota-se \(\eta_{\mu\nu}=\operatorname{diag}(1,-1,-1,-1)\), com
índices gregos de 0 a 3 e espaciais de 1 a 3. Em unidades \(c=1\),
uma perturbação plana é escrita como
\begin{equation}
 g_{\mu\nu}=\eta_{\mu\nu}
       +\operatorname{Re}\!\left[H_{\mu\nu}e^{-i\omega t+ikz}\right],
 \qquad q_\mu=(-\omega,0,0,k).
 \label{eq:fund-onda}
\end{equation}
As amplitudes \(H_{\mu\nu}\) são covariantes e podem ser complexas.
Para a métrica real, a frequência negativa contém o conjugado da
positiva. A dimensão do espaço de amplitudes em uma frequência não
deve ser duplicada por contabilizar artificialmente seno e cosseno como
novos estados de polarização. As identidades polinomiais verificadas
para amplitudes reais estendem-se por linearidade complexa.

A subida dos índices dá \(q^\mu=(-\omega,0,0,-k)\),
\(H^{0i}=-H_{0i}\), \(H^{00}=H_{00}\) e \(H^{ij}=H_{ij}\).
O traço e a reversão do traço são
\begin{equation}
 \begin{aligned}
 H&=H_{00}-H_{11}-H_{22}-H_{33},\\
 \bar H_{\mu\nu}&=H_{\mu\nu}-\tfrac12\eta_{\mu\nu}H,
 &\bar H&=-H.
 \end{aligned}
 \label{eq:fund-traco}
\end{equation}
Em quatro dimensões, a transformação inversa tem a mesma forma:
\(H_{\mu\nu}=\bar H_{\mu\nu}
-\eta_{\mu\nu}\bar H/2\). Essa distinção será essencial para
relacionar a notação do capítulo 4 do TG à entrada de seu apêndice.

Definem-se
\begin{equation}
 \Delta=\omega^2-k^2,\qquad \Sigma=\omega^2+k^2.
 \label{eq:fund-delta-sigma}
\end{equation}
As parametrizações serão demonstradas para \(\omega\ne0\) e
\(k\) real; nesse domínio, \(\Sigma>0\). Ondas massivas viajantes
pertencem a \(\Delta>0\). O caso nulo \(\Delta=0\) será avaliado
diretamente, e o limiar \(k=0\) não será removido por divisões
intermediárias. Para interpretar velocidades no sentido \(+z\),
escolhem-se \(\omega>0\) e \(k\geq0\). As provas algébricas de
posto também abrangem os demais sinais permitidos.

Em coordenadas \(x^0=ct\), o parâmetro temporal do cálculo é
\(\omega_{\mathrm{SI}}/c\). A dispersão operacional empregada no
projeto e suas escalas são
\begin{equation}
 \begin{aligned}
 \omega_{\mathrm{SI}}^2&=c^2k^2+
                  \left(\frac{m_gc^2}{\hbar}\right)^2,\\
 f_g&=\frac{m_gc^2}{h_{\mathrm P}},
 &\beta&=\frac{ck}{\omega_{\mathrm{SI}}}
        =\sqrt{1-(f_g/f)^2}.
 \end{aligned}
 \label{eq:fund-dispersao-si}
\end{equation}
A velocidade de grupo é \(c\beta\); a de fase, \(c/\beta\).
No limiar, a onda é espacialmente homogênea e não possui direção de
propagação privilegiada. Abaixo dele, \(k\) imaginário não descreve
uma onda viajante homogênea. A continuação algébrica de expressões não
basta para incluí-la na população de ondas viajantes de uma ORF.

\section{Curvatura elétrica sem hipótese de nulidade}

Adota-se a convenção
\begin{equation}
 R^\alpha{}_{\beta\gamma\delta}
 =\partial_\gamma\Gamma^\alpha{}_{\delta\beta}
  -\partial_\delta\Gamma^\alpha{}_{\gamma\beta}+O(H^2).
 \label{eq:fund-riemann-convencao}
\end{equation}
O símbolo de Christoffel linearizado,
\begin{equation}
 \Gamma^\alpha{}_{\mu\nu}
 =\tfrac12\eta^{\alpha\rho}
  (H_{\rho\mu,\nu}+H_{\rho\nu,\mu}-H_{\mu\nu,\rho}),
 \label{eq:fund-christoffel}
\end{equation}
fornece, com o fator de fase subentendido,
\begin{equation}
 \begin{aligned}
 R_{\alpha\beta\gamma\delta}^{(1)}
 =\tfrac12\big(&H_{\alpha\delta,\beta\gamma}
             +H_{\beta\gamma,\alpha\delta}\\
             &-H_{\alpha\gamma,\beta\delta}
             -H_{\beta\delta,\alpha\gamma}\big).
 \end{aligned}
 \label{eq:fund-riemann-linear}
\end{equation}
Cada derivada segunda da onda plana introduz \(-q_\mu q_\nu\).
Logo, as seis componentes independentes de
\(E_{ij}=R_{0i0j}\) são
\begin{equation}
 \begin{aligned}
 E_{11}&=\tfrac12\omega^2H_{11},
 &E_{22}&=\tfrac12\omega^2H_{22},\\
 E_{12}&=\tfrac12\omega^2H_{12},
 &E_{13}&=\tfrac12(\omega^2H_{13}+\omega kH_{01}),\\
 E_{23}&=\tfrac12(\omega^2H_{23}+\omega kH_{02}),\\
 E_{33}&=\tfrac12(\omega^2H_{33}+2\omega kH_{03}+k^2H_{00}).
 \end{aligned}
 \label{eq:fund-eletrica-geral}
\end{equation}
Não foram utilizados equação de campo, relação de dispersão, escolha
de gauge ou tetrada nula. Portanto, a existência de seis coordenadas
geométricas nessa matriz ainda não determina o número de amplitudes
físicas independentes.

O sinal de \eqref{eq:fund-riemann-linear} é oposto à saída do apêndice
Maple do TG. Na página impressa 44, este apresenta
\(R_{0x0x}^{\mathrm{TG}}=-\omega^2H_{11}/2\), enquanto
\eqref{eq:fund-eletrica-geral} fornece o sinal positivo. Com a
assinatura e a convenção aqui adotadas, a aceleração relativa tem
componentes espaciais euclidianas
\(a^i=c^2E_{ij}\xi^j\). Se uma matriz de força for definida por
\(a^i=-c^2S_{ij}\xi^j\), então \(S_{ij}=-E_{ij}\). Postos e
quocientes considerados adiante são invariantes sob essa inversão
global, mas as expressões de desvio geodésico não podem ser combinadas
com um sinal de curvatura incompatível.

As duas construções de Riemann, pela conexão e pela fórmula explícita,
permitem verificações independentes. Além das antissimetrias nos pares
e da simetria de troca de pares, a identidade algébrica de Bianchi é
\(R_{abcd}+R_{acdb}+R_{adbc}=0\). Em Fourier, a diferencial reduz-se a
\begin{equation}
 q_eR_{abcd}+q_cR_{abde}+q_dR_{abec}=0.
 \label{eq:fund-bianchi-fourier}
\end{equation}
Outra verificação é inserir
\(H^{\xi}_{\mu\nu}=q_\mu\xi_\nu+q_\nu\xi_\mu\), absorvendo
o fator imaginário na amplitude \(\xi\). Os quatro termos da
curvatura cancelam identicamente. Isso expressa a invariância da
curvatura linear sob uma perturbação de coordenadas; não estabelece uma
liberdade de gauge do campo massivo em fundo fixado.

% Seção produzida pela auditoria independente de ação, unidades e massa.
\section{Einstein linear e normalização histórica da massa}
\label{sec:fund-normalizacao}

A reprodução geométrica depende de \(\omega\), \(k\) e dos vínculos.
Relacioná-la a uma massa em unidades físicas exige ainda conferir o
coeficiente da equação de propagação. Nesta seção, \(\Box\) designa o
d'Alembertiano de segunda ordem. O símbolo \(\Box^2\) utilizado no TG
representa esse operador, não sua composição de quarta ordem.

Uma verificação independente, com dez componentes métricas e quatro
componentes de \(q_\mu\) livres, parte da conexão e contrai Ricci e
Einstein. Definindo \(B_\nu=q^\alpha\bar H_{\alpha\nu}\), obtém-se
\begin{equation}
 G_{\mu\nu}^{(1)}=-\frac12\Box\bar H_{\mu\nu}
 -\frac12(q_\mu B_\nu+q_\nu B_\mu)
 +\frac12\eta_{\mu\nu}q^\alpha B_\alpha,
 \label{eq:fund-einstein-geral}
\end{equation}
com \(\Box\mapsto-q^\alpha q_\alpha\) na amplitude de Fourier.
A contração \(q^\mu G_{\mu\nu}^{(1)}\) é identicamente nula.
Quando \(B_\nu=0\), resulta
\(G_{\mu\nu}^{(1)}=-\Box\bar H_{\mu\nu}/2\). A convenção oposta
de Riemann usada nas equações históricas inverte esse sinal:
\begin{equation}
 G_{\mu\nu}^{(1),\mathrm{hist}}=
       \frac12\Box\bar H_{\mu\nu}.
 \label{eq:fund-einstein-historico}
\end{equation}
A identidade foi conferida para as duas assinaturas de Minkowski.
Mudar a assinatura modifica a expressão de \(\Box\), mas não elimina
o fator \(1/2\).

As páginas impressas 23--24 do TG fornecem, nas equações (4.8) e
(4.9), o coeficiente linear
\(T_{\mu\nu}^{\mathrm{mass}}
=-m_{\mathrm{TG}}^2c^2\bar H_{\mu\nu}/(8\pi G\hbar^2)\)
e \(G_{\mu\nu}=8\pi GT_{\mu\nu}^{\mathrm{mass}}/c^4\) em vácuo
\autocite{tg2003}. Sua substituição literal em
\eqref{eq:fund-einstein-historico} produz
\begin{equation}
 \left(\Box+\frac{2m_{\mathrm{TG}}^2}{\hbar^2c^2}\right)
 \bar H_{\mu\nu}=0,
 \label{eq:fund-tg-literal}
\end{equation}
reproduzindo o fator dois das equações (4.17)--(4.19). Há, entretanto,
uma questão dimensional: se \(m_{\mathrm{TG}}\) é massa SI e
\(x^0=ct\), então
\begin{equation}
 \left[\frac{m_{\mathrm{TG}}^2}{\hbar^2c^2}\right]=L^{-6}T^4,
 \qquad
 \left[\left(\frac{m_gc}{\hbar}\right)^2\right]=L^{-2}.
 \label{eq:fund-dimensoes-massa}
\end{equation}
Falta um fator \(c^4\) no primeiro coeficiente relativamente à dimensão
exigida. Reinterpretar o símbolo como massa previamente reescalada
exigiria explicitar essa definição e traduzir também a velocidade
impressa após (4.22). Usar \(c=1\) oculta a questão dimensional,
sem demonstrar a equivalência das identificações de massa.

O artigo de 2004 adota assinatura \((-+++)\) e escreve \(c^6\) no
numerador de \(T^{\mathrm{mass}}\), nas equações (4)--(5), restaurando
a dimensão de densidade de energia \autocite{depaula2004}. Sua equação
(6) usa \(G=-8\pi GT^{\mathrm{mass}}/c^4\). Denotando por
\(m_a\) o parâmetro dessa fonte e conservando o sinal de Ricci da
equação (10), a cadeia literal resulta em
\begin{equation}
 \left(\Box-2\frac{m_a^2c^2}{\hbar^2}\right)\bar H_{\mu\nu}=0.
 \label{eq:fund-artigo-literal}
\end{equation}
Já a equação (12) imprime o coeficiente \(-m_a^2c^2/\hbar^2\).
Portanto, as expressões consultadas apresentam uma discrepância de
fator dois se o parâmetro de massa for o mesmo. A mudança de assinatura
explica o sinal relativo ao TG, mas não esse fator. Reduzir o
coeficiente fonte pela metade ou definir
\(m_{\mathrm{disp}}^2=2m_a^2\) são escolhas de normalização distintas;
não se devem aplicar ambas como compensações silenciosas.

Este é um diagnóstico algébrico das equações consultadas, não uma
errata editorial confirmada ou uma reconstrução da ação não linear.
A dissertação adota a \emph{massa de dispersão} \(m_g\), definida por
\eqref{eq:fund-dispersao-si}, e \(\mu=m_gc/\hbar\).
Esse coeficiente operacional também aparece na formulação linear de
Visser \autocite{visser1998}. Sua identificação com um parâmetro de
outra normalização deve ser demonstrada em cada comparação. As tabelas
em \(\mathrm{eV}/c^2\) e a resposta futura de PTA usarão a definição
operacional; a reprodução das expressões de curvatura permanecerá em
\(\omega\) e \(k\), preservando as fontes históricas.


\section{Vínculos do modelo histórico e seis amplitudes}

O vínculo linear associado à equação (4.14) do TG é
\begin{equation}
 F_\nu\equiv q^\mu H_{\mu\nu}-\tfrac12q_\nu H
          =q^\mu\bar H_{\mu\nu}=0.
 \label{eq:fund-visser-vinculo}
\end{equation}
No contexto histórico, ele decorre das equações e das condições de
conservação adotadas \autocite{visser1998,tg2003}. A semelhança com uma
condição harmônica da relatividade geral não autoriza eliminar modos
massivos por gauge. De fato, para a perturbação de coordenadas
anterior, obtém-se \(F_\nu(H^\xi)=\Delta\xi_\nu\). Quando
\(\Delta\ne0\), uma transformação não trivial desse tipo em geral
não preserva o vínculo.

As quatro equações podem ser exibidas como
\begin{equation}
 \begin{aligned}
 \omega(H_{00}+H_{11}+H_{22}+H_{33})+2kH_{03}&=0,\\
 \omega H_{01}+kH_{13}&=0,\\
 \omega H_{02}+kH_{23}&=0,\\
 2\omega H_{03}+k(H_{00}-H_{11}-H_{22}+H_{33})&=0.
 \end{aligned}
 \label{eq:fund-visser-equacoes}
\end{equation}
Escolhem-se as amplitudes livres
\begin{equation}
 (A,D,B,C,V_x,V_y)
 =(H_{00},H_{33},H_{22},H_{12},H_{13},H_{23}).
 \label{eq:fund-amplitudes-livres}
\end{equation}
Resolver \eqref{eq:fund-visser-equacoes} fornece
\begin{equation}
 \begin{aligned}
 H_{01}&=-\frac{k}{\omega}V_x,
 &H_{02}&=-\frac{k}{\omega}V_y,\\
 H_{03}&=-\frac{\omega k}{\Sigma}(A+D),
 &H_{11}&=-B-\frac{\Delta}{\Sigma}(A+D).
 \end{aligned}
 \label{eq:fund-visser-solucao}
\end{equation}
Não houve divisão por \(k\) nem por \(\Delta\). A parametrização
é regular tanto no limiar quanto no ponto nulo do domínio.

Para demonstrar completude, considera-se o jacobiano de
\((F_0,F_1,F_2,F_3)\) nas colunas
\((H_{01},H_{02},H_{03},H_{11})\):
\begin{equation}
 M_V=\begin{pmatrix}
 0&0&-k&-\omega/2\\
 -\omega&0&0&0\\
 0&-\omega&0&0\\
 0&0&-\omega&k/2
 \end{pmatrix},
 \qquad \det M_V=-\frac{\omega^2\Sigma}{2}.
 \label{eq:fund-menor-visser}
\end{equation}
Como esse menor é não nulo e há quatro linhas, o posto dos vínculos é
exatamente quatro para todo \(\omega\ne0\), \(k\) real. Das dez
componentes simétricas restam seis amplitudes. Essa conclusão não
depende de uma escolha numérica genérica feita por um sistema
computacional.

Escrevendo \(U=A+D\), sua matriz de marés é
\begin{equation}
 E_V=\frac12\begin{pmatrix}
 -\omega^2(B+\Delta U/\Sigma)&\omega^2C&\Delta V_x\\
 \omega^2C&\omega^2B&\Delta V_y\\
 \Delta V_x&\Delta V_y&\Delta(\omega^2D-k^2A)/\Sigma
 \end{pmatrix}.
 \label{eq:fund-eletrica-visser}
\end{equation}
Definem-se coordenadas de deformação com normalização explícita:
\begin{equation}
 \mathbf p=(p_b,p_l,p_+,p_\times,p_x,p_y)
 =(E_{11}+E_{22},E_{33},E_{11}-E_{22},E_{12},E_{13},E_{23}).
 \label{eq:fund-coordenadas-mares}
\end{equation}
Em particular,
\begin{equation}
 \begin{aligned}
 p_b&=-\frac{\omega^2\Delta}{2\Sigma}(A+D),
 &p_l&=\frac{\Delta}{2\Sigma}(\omega^2D-k^2A),\\
 p_+&=p_b-\omega^2B,
 &p_\times&=\frac{\omega^2C}{2},\\
 p_x&=\frac{\Delta V_x}{2},
 &p_y&=\frac{\Delta V_y}{2}.
 \end{aligned}
 \label{eq:fund-coordenadas-visser}
\end{equation}
O bloco de \((p_b,p_l)\) em \((A,D)\) tem determinante
\(-\omega^2\Delta^2/(4\Sigma)\). Subtrair a linha de \(p_b\) da
de \(p_+\) separa as demais quatro amplitudes e resulta em
\begin{equation}
 \det\frac{\partial\mathbf p}{\partial(A,D,B,C,V_x,V_y)}
       =\frac{\omega^6\Delta^4}{32\Sigma}.
 \label{eq:fund-posto-visser}
\end{equation}
O mapa tem, portanto, posto seis quando \(\Delta\ne0\). Isso é
uma contagem cinemática do modelo linear considerado, e não uma prova
de seis modos dinamicamente saudáveis.

\section{Fierz--Pauli e o vínculo entre deformações escalares}

Uma normalização compatível com a convenção de curvatura é
\begin{equation}
 G_{\mu\nu}^{(1)}
    -\frac{\mu^2}{2}(H_{\mu\nu}-\eta_{\mu\nu}H)=0,
 \qquad \mu\ne0.
 \label{eq:fund-fp-campo}
\end{equation}
A identidade de Bianchi implica
\(\partial^\mu H_{\mu\nu}=\partial_\nu H\). Substituí-la em
\(R^{(1)}=\partial_\mu\partial_\nu H^{\mu\nu}-\Box H\)
fornece \(R^{(1)}=0\). O traço de
\eqref{eq:fund-fp-campo} é então \(3\mu^2H/2=0\), de modo que
\begin{equation}
 H=0,\qquad q^\mu H_{\mu\nu}=0,
 \qquad (\Box+\mu^2)H_{\mu\nu}=0.
 \label{eq:fund-fp-vinculos}
\end{equation}
A última equação dá \(\Delta=\mu^2\). Dividir por \(\mu^2\)
foi essencial; essa derivação dos vínculos não pode ser repetida em
massa exatamente zero como se nada mudasse.

Na parametrização anterior, o traço impõe
\(A=k^2D/\omega^2\), e as componentes eliminadas tornam-se
\begin{equation}
 H_{03}=-\frac{k}{\omega}D,
 \qquad H_{11}=-B-\frac{\Delta}{\omega^2}D,
 \qquad H_{01,02}=-\frac{k}{\omega}V_{x,y}.
 \label{eq:fund-fp-solucao}
\end{equation}
O jacobiano dos quatro vínculos de divergência seguidos do traço,
nas colunas \((H_{00},H_{01},H_{02},H_{03},H_{11})\), contém o menor
\begin{equation}
 M_F=\begin{pmatrix}
 -\omega&0&0&-k&0\\
 0&-\omega&0&0&0\\
 0&0&-\omega&0&0\\
 0&0&0&-\omega&0\\
 1&0&0&0&-1
 \end{pmatrix},
 \qquad\det M_F=-\omega^4.
 \label{eq:fund-menor-fp}
\end{equation}
Restam exatamente cinco amplitudes no domínio considerado. A matriz
elétrica é
\begin{equation}
 E_F=\frac12\begin{pmatrix}
 -\omega^2B-\Delta D&\omega^2C&\Delta V_x\\
 \omega^2C&\omega^2B&\Delta V_y\\
 \Delta V_x&\Delta V_y&\Delta^2D/\omega^2
 \end{pmatrix}.
 \label{eq:fund-eletrica-fp}
\end{equation}
O menor formado pelas linhas \((p_b,p_+,p_\times,p_x,p_y)\) e
colunas \((B,C,V_x,V_y,D)\) vale
\begin{equation}
 \det J_F^{(5)}=\frac{\omega^4\Delta^3}{16}.
 \label{eq:fund-posto-fp}
\end{equation}
Assim, para \(\Delta\ne0\), o posto observável é cinco. As duas
formas escalares não podem constituir amplitudes independentes:
\begin{equation}
 p_b=-\frac{\Delta D}{2},\qquad
 p_l=\frac{\Delta^2D}{2\omega^2},\qquad
 p_l+\frac{\Delta}{\omega^2}p_b=0.
 \label{eq:fund-identidade-escalar}
\end{equation}
A identidade permanece válida quando \(p_b=0\), situação em que
uma razão não seria definida. Quando \(p_b\ne0\), ela equivale a
\begin{equation}
 \boxed{\frac{p_l}{p_b}=-(1-\beta^2)=-\left(\frac{f_g}{f}\right)^2.}
 \label{eq:fund-razao-escalar}
\end{equation}
O modo breathing foi normalizado como a soma das duas componentes
transversais. Usar apenas uma delas para uma deformação transversal
isotrópica introduz um fator dois na razão, sem alterar o vínculo
físico. Para isolar a helicidade zero na base adotada, impõem-se também
\(C=V_x=V_y=0\) e \(E_{11}=E_{22}\), isto é,
\(B=-\Delta D/(2\omega^2)\). Essa única amplitude pode produzir
simultaneamente deformações transversal e longitudinal
\autocite{liang2021,bernardo2024}.

\section{Limite nulo, limiar e não uniformidade}

Avaliar diretamente \(\Delta=0\), com \(\omega\ne0\), nas
parametrizações regulares fornece, nos dois casos,
\begin{equation}
 E_{\mathrm{nulo}}=\frac{\omega^2}{2}
 \begin{pmatrix}-B&C&0\\ C&B&0\\0&0&0\end{pmatrix}.
 \label{eq:fund-eletrica-nula}
\end{equation}
Seu posto como mapa de amplitudes é dois. A dimensão do núcleo é
quatro no espaço de seis amplitudes histórico e três no espaço de
cinco amplitudes com os vínculos de Fierz--Pauli formalmente mantidos.
Essa afirmação é algébrica e não pressupõe que a equação massiva
derive os mesmos vínculos na teoria sem massa.

No espaço harmônico nulo, todas as quatro amplitudes
\(H^\xi_{\mu\nu}=q_\mu\xi_\nu+q_\nu\xi_\mu\) satisfazem o
vínculo e têm curvatura zero. O mapa \(\xi\mapsto H^\xi\) é
injetivo para \(\omega\ne0\): a componente 00 determina
\(\xi_0\), e as 0i determinam as demais. No subespaço com traço
nulo, exige-se ainda \(q^\mu\xi_\mu=0\), deixando três
direções. Essas dimensões coincidem com os núcleos acima e identificam
as direções sem curvatura nesse caso nulo específico.

No limiar \(k=0\), os menores massivos continuam não nulos. O
modelo histórico conserva posto seis e Fierz--Pauli, cinco. Neste,
\(H_{00}=H_{0i}=0\), e a matriz espacial tem traço zero. A relação
escalar torna-se \(p_l=-p_b\). Como não há direção de propagação
preferencial, rótulos de helicidade associados a um eixo escolhido não
devem ser confundidos com uma mudança na contagem de estados.

O ponto \(q_\mu=0\) exige avaliação separada: as derivadas e a
curvatura desaparecem. Os quatro vínculos históricos tornam-se
nulos; a imposição formal de traço zero conserva uma equação. Isso
não descreve uma onda massiva no vácuo, pois
\(\Delta=\mu^2>0\) não pode ser satisfeita nesse ponto. Campos
estáticos com \(\omega=0\) também estão fora da parametrização
utilizada e não devem ser obtidos por divisões por zero.

O desaparecimento de componentes não tensoriais em
\eqref{eq:fund-eletrica-nula} refere-se a famílias com amplitudes
métricas limitadas. Uma parametrização alternativa que usa todas as
componentes espaciais livres exige
\begin{equation}
 H_{00}=-H_{33}-\frac{\Sigma}{\Delta}(H_{11}+H_{22}).
 \label{eq:fund-parametrizacao-singular}
\end{equation}
Manter a soma transversal fixa pode, portanto, levar a uma métrica
ilimitada quando \(\Delta\to0\). Não é legítimo trocar o limite
e a escolha da família de amplitudes sem examinar essa dependência.

Uma ilustração explícita em Fierz--Pauli é
\(D=-2\omega^2/\Delta\), \(B=1\) e \(C=V_x=V_y=0\):
\begin{equation}
 H_{\mathrm{esc}}=\begin{pmatrix}
 -2k^2/\Delta&0&0&2\omega k/\Delta\\
 0&1&0&0\\
 0&0&1&0\\
 2\omega k/\Delta&0&0&-2\omega^2/\Delta
 \end{pmatrix},
 \label{eq:fund-familia-singular}
\end{equation}
para a qual
\begin{equation}
 E_{\mathrm{esc}}=\operatorname{diag}
       (\omega^2/2,\omega^2/2,-\Delta).
 \label{eq:fund-familia-curvatura}
\end{equation}
A deformação breathing não desaparece, mas a métrica não permanece
limitada. Para uma amplitude global fixa, não se pode prolongar essa
família até \(\Delta=0\) mantendo o regime linear. Rescalá-la para
restabelecer pequenas perturbações também altera o limite de
curvatura. O exemplo demonstra não uniformidade matemática, não a
continuidade com a relatividade geral, uma previsão de emissão ou uma
derivação da descontinuidade vDVZ. Essas questões exigem acoplamento a
fontes e análise dinâmica.

\section{Rastreamento das projeções do TG}

\subsection{Amplitude métrica, tetradas e sinal de curvatura}

A equação (4.21) do TG chama \(A_{\mu\nu}\) à amplitude de
\(\bar h_{\mu\nu}\), mas o apêndice insere o mesmo nome diretamente
em \(g_{\mu\nu}=\eta_{\mu\nu}+A_{\mu\nu}e^{i\phi}\).
A reprodução abaixo usa a amplitude \emph{métrica} do apêndice, com
\(A_{00}=A\), \(A_{33}=D\), \(A_{22}=B\), \(A_{12}=C\),
\(A_{13}=V_x\) e \(A_{23}=V_y\). Se o traço for não nulo, esses
coeficientes não podem ser identificados diretamente com os de
\(\bar h\) sem aplicar \eqref{eq:fund-traco}.

A tetrada simétrica do apêndice pode ser escrita como
\begin{equation}
 \begin{aligned}
 l^\mu&=\frac{(1,0,0,1)}{\sqrt2},
 &n^\mu&=\frac{(1,0,0,-1)}{\sqrt2},\\
 m^\mu&=\frac{(0,1,i,0)}{\sqrt2},
 &\bar m^\mu&=\frac{(0,1,-i,0)}{\sqrt2}.
 \end{aligned}
 \label{eq:fund-tetrada}
\end{equation}
Em nossa assinatura, \(l\cdot n=1\) e
\(m\cdot\bar m=-1\). Os sinais da relação (5.19) do corpo do TG
devem ser traduzidos para essa convenção. Seus vetores longitudinais
nas equações (5.12)--(5.13) são \(\sqrt2\,l\) e
\(n/\sqrt2\); logo, a contração com dois vetores \(n\) muda por
um fator dois. Mantendo a mesma orientação complexa, a projeção
\(\Psi_4\) na tetrada do corpo vale metade daquela do apêndice.
Essa transformação é um boost de tetrada, não uma nova polarização.

\subsection{Expressões reduzidas e contração independente}

As transcrições literais das equações (5.58)--(5.61) simplificam para
\begin{equation}
 \begin{aligned}
 \Psi_2^{\mathrm{TG}}&=
    \frac{\Delta}{12\Sigma}(\omega^2D-k^2A),\\
 \Psi_3^{\mathrm{TG}}&=
    \frac{\Delta(\omega+k)}{8\omega}(V_x+iV_y),\\
 \Psi_4^{\mathrm{TG}}&=
    \frac{(\omega+k)^2}{8}
        \left[-2B-\frac{\Delta}{\Sigma}(A+D)+2iC\right],\\
 \Phi_{22}^{\mathrm{TG}}&=
    -\frac{\Delta(\omega+k)^2}{8\Sigma}(A+D).
 \end{aligned}
 \label{eq:fund-tg-projecoes}
\end{equation}
Essas identidades foram confrontadas com contrações do Riemann
construído independentemente. Escrevendo \(R(a,b,c,d)\) para a
contração com quatro vetores, a convenção adotada fornece
\begin{equation}
 \begin{aligned}
 \Psi_2^{\mathrm{TG}}&=R(n,l,n,l)/6,
 &\Psi_3^{\mathrm{TG}}&=R(n,l,n,m)/2,\\
 \Psi_4^{\mathrm{TG}}&=R(n,m,n,m),
 &\Phi_{22}^{\mathrm{TG}}&=R(n,\bar m,n,m).
 \end{aligned}
 \label{eq:fund-tg-contracoes}
\end{equation}
O sinal acompanha \(R_{\mathrm{TG}}=-R\). A orientação de \(m\)
foi escolhida para reproduzir o sinal positivo da parte imaginária de
\(\Psi_4\) impresso no TG. Conjugar a direção transversal sem
acompanhar a convenção trocaria esse sinal.

Reproduzir essas expressões não demonstra que toda sua interpretação
nula permaneça válida em propagação massiva. As equações (5.35) e
(5.36) do TG relacionam as projeções a escalares de Weyl com correções
de ordem quase nula explicitamente omitidas. O apêndice efetivamente
contrai Riemann. Para \(k\ne\omega\), seus nomes
\(\Psi_2^{\mathrm{TG}}\) e \(\Psi_3^{\mathrm{TG}}\) não devem
ser identificados automaticamente com os escalares exatos de Weyl.
Na contração \(nmnm\), os termos de Ricci da decomposição de Riemann
se anulam por nulidade e ortogonalidade dos vetores: esse caso
particular não sofre a mesma ambiguidade.

O grupo E(2) preserva um vetor de onda nulo. Para um vetor temporal
massivo, o estabilizador de Lorentz é SO(3). Uma tetrada nula continua
sendo uma base legítima, mas isso não transforma uma classificação
quase nula em classificação exata da teoria massiva. A matriz elétrica
é utilizada aqui para manter a ligação explícita com o observador.

\subsection{Benchmark tensorial de uma conversão quase nula}

Para \(A=D=B=V_x=V_y=0\) e \(C=1\) real, com fase comum,
\eqref{eq:fund-tg-projecoes} dá
\(\operatorname{Im}\Psi_4^{\mathrm{TG}}=(\omega+k)^2/4\),
enquanto a componente exata é \(E_{12}=\omega^2/2\).
A conversão calibrada na onda nula utiliza
\(\operatorname{Im}\Psi_4^{\mathrm{TG}}/2\). Seu quociente com
o valor exato é
\begin{equation}
 \mathcal R_T(\beta)
 =\frac{\operatorname{Im}\Psi_4^{\mathrm{TG}}/2}{E_{12}}
 =\frac{(1+\beta)^2}{4}.
 \label{eq:fund-erro-tensorial}
\end{equation}
O resultado tende a um para \(\beta\to1\) e vale \(1/4\) no
limiar, produzindo erro relativo de 75\% nessa conversão. Trata-se de
uma projeção tensorial e de sua tradução para curvatura, não de uma
amplitude de polarização adicional ou de um erro medido em dados PTA.

Com \(x=(f_g/f)^2\), uma forma numericamente estável do erro é
\begin{equation}
 1-\mathcal R_T
 =\frac{x(3+\sqrt{1-x})}{4(1+\sqrt{1-x})}
 =\frac{x}{2}+O(x^2).
 \label{eq:fund-erro-estavel}
\end{equation}
Essa expressão evita subtrair dois números quase iguais quando
\(x\) é pequeno. Uma implementação pode reproduzir exatamente a
álgebra do TG e ainda precisar dessa auditoria para converter as
projeções em componentes físicas fora do domínio nulo.

\section{Confronto com Hyun e escalas numéricas}

As equações (3.29)--(3.30) de Hyun, Kim e Lee, na página 14 do PDF
arXiv consultado, empregam as componentes livres
\((h_{tt},h_{zz},h_{yy},h_{xy},h_{xz},h_{yz})\)
\autocite{hyun2019}. Após traduzir assinatura, sinal e ordenação, o
vetor de suas seis amplitudes coincide com
\eqref{eq:fund-coordenadas-visser}. Por exemplo, sua componente
longitudinal pode ser escrita como
\begin{equation}
 \frac{\omega^2\Delta}{2\Sigma}(A+D)-\frac{\Delta A}{2}
 =\frac{\Delta}{2\Sigma}(\omega^2D-k^2A)=p_l.
 \label{eq:fund-hyun-longitudinal}
\end{equation}
A transcrição independente das seis componentes foi verificada
simbolicamente. A expressão equivalente com componentes espaciais
livres, equação (3.27) da referência, requer atenção à singularidade de
parametrização discutida anteriormente. A concordância é um benchmark
de reprodução de um resultado conhecido, sem reivindicação de
originalidade.

Para dimensionar a relevância das escalas, utiliza-se como benchmark
\(m_gc^2=1{,}92\times10^{-23}\,\mathrm{eV}\), valor associado ao
limite de dispersão reportado em GWTC-4 \autocite{lvk2026}. Não se
trata de massa detectada ou limite universal entre modelos. Com
\(h_{\mathrm P}=4{,}135667696\ldots\times10^{-15}\,
\mathrm{eV\,s}\), resulta
\(f_g=4{,}642539345\ldots\times10^{-9}\,\mathrm{Hz}\).
As avaliações do script preliminar fornecem:

\begin{center}
\begingroup\small\singlespacing
\setlength{\tabcolsep}{6pt}
\begin{tabular}{@{}rccc@{}}
\toprule
\(f\) (Hz) & \(x=(f_g/f)^2\) & \(\beta\) & \(1-\mathcal R_T\)\\
\midrule
\(10^{-9}\) & \(21{,}5532\) & não viajante & não aplicável\\
\(10^{-8}\) & \(0{,}215532\) & \(0{,}885702\) & \(0{,}111032\)\\
\(10^{-7}\) & \(0{,}00215532\) & \(0{,}998922\) & \(0{,}00107795\)\\
\(10^{-4}\) & \(2{,}15532\times10^{-9}\) & \(\simeq1\) & \(1{,}07766\times10^{-9}\)\\
\(100\) & \(2{,}15532\times10^{-21}\) & \(\simeq1\) & \(1{,}07766\times10^{-21}\)\\
\bottomrule
\end{tabular}
\endgroup
\end{center}

Para a escolha histórica do artigo de 2004,
\(m_gc^2=4{,}4\times10^{-22}\,\mathrm{eV}\) e
\(f=1{,}1\times10^{-7}\,\mathrm{Hz}\), obtêm-se
\(x\simeq0{,}935468\), \(\beta\simeq0{,}254032\) e
\(1-\mathcal R_T\simeq0{,}606851\)
\autocite{depaula2004}. O regime está longe da onda nula nessa
conversão operacional de unidades. Isso não equivale a declarar um
ajuste observacional histórico incorreto: a comparação verifica o
domínio de uma aproximação cinemática. A relação entre os símbolos de
massa das equações históricas e a massa operacional deve seguir a
auditoria de normalização deste capítulo.

\section{Evidência computacional e alcance dos resultados}

A implementação simbólica integrada a esta versão registrou 19 testes
aprovados, sem falhas, com SymPy 1.14.0: doze de polarizações e sete
da auditoria independente de Einstein, dinâmica linear e normalização. Os testes comparam as 256
componentes de Riemann obtidas por duas rotas, verificam simetrias e
as duas identidades de Bianchi, e anulam a curvatura de perturbações de
coordenadas. Também confrontam as parametrizações com a solução direta
dos sistemas de vínculos, demonstram os menores, avaliam os casos
nulo e de limiar, verificam o vínculo escalar e sua família não
uniforme, e reproduzem as expressões de Hyun e do TG.

Os cálculos simbólicos gerais ampliam as checagens racionais
preliminares em \((\omega,k)=(1,0),(5,3),(13,12),(1,1)\). Esses
exemplos continuam úteis como regressão, mas a prova dos postos é dada
pelos menores explícitos e seus domínios. As escalas físicas foram
avaliadas com precisão decimal ampliada, evitando que o arredondamento
oculte correções pequenas. O módulo \path{src/polarizacoes/}, os testes em \path{tests/} e o
script \path{scripts/reproduzir_tg.py} acompanham esta versão. O registro
\path{results/C04/validacao.json} contém os testes executados, as
expressões simbólicas e as dependências fixadas. A opção \texttt{-{}-check}
reexecuta a validação sem substituir o registro versionado.

O alcance desses resultados é cinemático. Não foram calculadas neste
capítulo ORFs, evidências bayesianas, cobertura de intervalos, excitação
por fontes ou estabilidade não linear. A relação local entre
\(p_l\) e \(p_b\) não é uma função de resposta de PTA. Em
particular, a revisão v3 de Liang--Trodden exige incluir componentes
temporais da métrica e o movimento dos observadores na medição da
frequência \autocite{liang2021}. A passagem da curvatura local para
resíduos de temporização será desenvolvida e validada em C05, e a
população escalar vinculada será examinada em C10. Nenhuma dessas
validações estatísticas futuras é substituída pelos 19 testes aqui
registrados.
